% ==========================================
% ANEXOS — ARCHIVO AGREGADOR
% ==========================================
%
% Este archivo es el punto de entrada de todos los anexos.
% No pongas \chapter{} aquí directamente (a menos que sea
% un anexo general). En su lugar, crea un anexo por capítulo
% y agrégalo con \input{}.
%
% TIP: Cada capítulo puede tener su propio anexo en:
%      capitulos/capitulo_X/anexos.tex
%
% TIP: Los anexos se numeran con letras (A, B, C...)
%      automáticamente según el orden de los \input{}.
% ==========================================

\appendix
\renewcommand{\appendixname}{Anexo}
\renewcommand{\appendixpagename}{Anexos}
\renewcommand{\appendixtocname}{Anexos}
\addappheadtotoc
\appendixpage

% --- Agrega aquí los anexos de cada capítulo ---
% Ejemplo:
% % ==========================================
% ANEXOS — ARCHIVO AGREGADOR
% ==========================================
%
% Este archivo es el punto de entrada de todos los anexos.
% No pongas \chapter{} aquí directamente (a menos que sea
% un anexo general). En su lugar, crea un anexo por capítulo
% y agrégalo con \input{}.
%
% TIP: Cada capítulo puede tener su propio anexo en:
%      capitulos/capitulo_X/anexos.tex
%
% TIP: Los anexos se numeran con letras (A, B, C...)
%      automáticamente según el orden de los \input{}.
% ==========================================

\appendix
\renewcommand{\appendixname}{Anexo}
\renewcommand{\appendixpagename}{Anexos}
\renewcommand{\appendixtocname}{Anexos}
\addappheadtotoc
\appendixpage

% --- Agrega aquí los anexos de cada capítulo ---
% Ejemplo:
% % ==========================================
% ANEXOS — ARCHIVO AGREGADOR
% ==========================================
%
% Este archivo es el punto de entrada de todos los anexos.
% No pongas \chapter{} aquí directamente (a menos que sea
% un anexo general). En su lugar, crea un anexo por capítulo
% y agrégalo con \input{}.
%
% TIP: Cada capítulo puede tener su propio anexo en:
%      capitulos/capitulo_X/anexos.tex
%
% TIP: Los anexos se numeran con letras (A, B, C...)
%      automáticamente según el orden de los \input{}.
% ==========================================

\appendix
\renewcommand{\appendixname}{Anexo}
\renewcommand{\appendixpagename}{Anexos}
\renewcommand{\appendixtocname}{Anexos}
\addappheadtotoc
\appendixpage

% --- Agrega aquí los anexos de cada capítulo ---
% Ejemplo:
% % ==========================================
% ANEXOS — ARCHIVO AGREGADOR
% ==========================================
%
% Este archivo es el punto de entrada de todos los anexos.
% No pongas \chapter{} aquí directamente (a menos que sea
% un anexo general). En su lugar, crea un anexo por capítulo
% y agrégalo con \input{}.
%
% TIP: Cada capítulo puede tener su propio anexo en:
%      capitulos/capitulo_X/anexos.tex
%
% TIP: Los anexos se numeran con letras (A, B, C...)
%      automáticamente según el orden de los \input{}.
% ==========================================

\appendix
\renewcommand{\appendixname}{Anexo}
\renewcommand{\appendixpagename}{Anexos}
\renewcommand{\appendixtocname}{Anexos}
\addappheadtotoc
\appendixpage

% --- Agrega aquí los anexos de cada capítulo ---
% Ejemplo:
% \input{capitulos/capitulo_1/anexos}
% \input{capitulos/capitulo_2/anexos}

% ==========================================
% REFERENCIA: Cómo incluir código fuente
% ==========================================
%
% Para proyectos pequeños, puedes poner los anexos
% directamente aquí. Para proyectos grandes, crea un
% archivo por capítulo (ver sección "Documentos grandes"
% en el README.md).
%
% A continuación se muestran las formas más comunes
% de incluir código con listings.
%
% ==========================================
%
% --- FORMA 1: Desde archivo externo (recomendada) ---
% Crea una carpeta codigos/ o programa/ en la raíz.
%
% \chapter{Código fuente}
% \label{anexo:codigo}
%
% \lstinputlisting[language=C, style=mystyle,
%   caption={Descripción del código. \propia{}},
%   label={lst:etiqueta}]{codigos/mi_archivo.c}
%
%
% --- FORMA 2: Con verificación de existencia ---
% Si el archivo puede no estar presente, usa \IfFileExists
% para evitar errores de compilación.
%
% \IfFileExists{programa/mi_programa.c}{
%   \lstinputlisting[language=C, style=mystyle,
%     caption={Código completo. \propia{}},
%     label={lst:codigo}]{programa/mi_programa.c}
% }{
%   \noindent\textbf{Nota:} El archivo aún no está disponible.
% }
%
%
% --- FORMA 3: Código inline (fragmentos pequeños) ---
% \begin{lstlisting}[language=Python, style=mystyle,
%   caption={Fragmento de ejemplo},
%   label={lst:ejemplo}]
% def fibonacci(n):
%     if n <= 1:
%         return n
%     return fibonacci(n - 1) + fibonacci(n - 2)
% \end{lstlisting}
%
%
% --- FORMA 4: Makefile ---
% \lstinputlisting[language=make, style=mystyle,
%   caption={Makefile del proyecto. \propia{}},
%   label={lst:makefile}]{codigos/Makefile}
%
%
% --- FORMA 5: Archivos de cabecera (.h) ---
% \lstinputlisting[language=C, style=mystyle,
%   caption={Archivo de cabecera. \propia{}},
%   label={lst:header}]{programa/common.h}
%
%
% ==========================================
% BUENAS PRÁCTICAS
% ==========================================
%
% 1. Usa \lstinputlisting para código que compila.
%    Así el PDF siempre refleja el código real.
%
% 2. Usa \IfFileExists para que la compilación no
%    falle si falta un archivo.
%
% 3. Pon siempre caption y label para referenciar
%    los listados desde el texto:
%    "En el Listado~\ref{lst:etiqueta} se muestra..."
%
% 4. Usa \propia{} al final del caption para agregar
%    "Fuente: Elaboración propia."
%
% 5. Para proyectos con varios capítulos, crea un
%    anexos.tex por capítulo y agrégalo arriba con
%    \input{}.

% % ==========================================
% ANEXOS — ARCHIVO AGREGADOR
% ==========================================
%
% Este archivo es el punto de entrada de todos los anexos.
% No pongas \chapter{} aquí directamente (a menos que sea
% un anexo general). En su lugar, crea un anexo por capítulo
% y agrégalo con \input{}.
%
% TIP: Cada capítulo puede tener su propio anexo en:
%      capitulos/capitulo_X/anexos.tex
%
% TIP: Los anexos se numeran con letras (A, B, C...)
%      automáticamente según el orden de los \input{}.
% ==========================================

\appendix
\renewcommand{\appendixname}{Anexo}
\renewcommand{\appendixpagename}{Anexos}
\renewcommand{\appendixtocname}{Anexos}
\addappheadtotoc
\appendixpage

% --- Agrega aquí los anexos de cada capítulo ---
% Ejemplo:
% \input{capitulos/capitulo_1/anexos}
% \input{capitulos/capitulo_2/anexos}

% ==========================================
% REFERENCIA: Cómo incluir código fuente
% ==========================================
%
% Para proyectos pequeños, puedes poner los anexos
% directamente aquí. Para proyectos grandes, crea un
% archivo por capítulo (ver sección "Documentos grandes"
% en el README.md).
%
% A continuación se muestran las formas más comunes
% de incluir código con listings.
%
% ==========================================
%
% --- FORMA 1: Desde archivo externo (recomendada) ---
% Crea una carpeta codigos/ o programa/ en la raíz.
%
% \chapter{Código fuente}
% \label{anexo:codigo}
%
% \lstinputlisting[language=C, style=mystyle,
%   caption={Descripción del código. \propia{}},
%   label={lst:etiqueta}]{codigos/mi_archivo.c}
%
%
% --- FORMA 2: Con verificación de existencia ---
% Si el archivo puede no estar presente, usa \IfFileExists
% para evitar errores de compilación.
%
% \IfFileExists{programa/mi_programa.c}{
%   \lstinputlisting[language=C, style=mystyle,
%     caption={Código completo. \propia{}},
%     label={lst:codigo}]{programa/mi_programa.c}
% }{
%   \noindent\textbf{Nota:} El archivo aún no está disponible.
% }
%
%
% --- FORMA 3: Código inline (fragmentos pequeños) ---
% \begin{lstlisting}[language=Python, style=mystyle,
%   caption={Fragmento de ejemplo},
%   label={lst:ejemplo}]
% def fibonacci(n):
%     if n <= 1:
%         return n
%     return fibonacci(n - 1) + fibonacci(n - 2)
% \end{lstlisting}
%
%
% --- FORMA 4: Makefile ---
% \lstinputlisting[language=make, style=mystyle,
%   caption={Makefile del proyecto. \propia{}},
%   label={lst:makefile}]{codigos/Makefile}
%
%
% --- FORMA 5: Archivos de cabecera (.h) ---
% \lstinputlisting[language=C, style=mystyle,
%   caption={Archivo de cabecera. \propia{}},
%   label={lst:header}]{programa/common.h}
%
%
% ==========================================
% BUENAS PRÁCTICAS
% ==========================================
%
% 1. Usa \lstinputlisting para código que compila.
%    Así el PDF siempre refleja el código real.
%
% 2. Usa \IfFileExists para que la compilación no
%    falle si falta un archivo.
%
% 3. Pon siempre caption y label para referenciar
%    los listados desde el texto:
%    "En el Listado~\ref{lst:etiqueta} se muestra..."
%
% 4. Usa \propia{} al final del caption para agregar
%    "Fuente: Elaboración propia."
%
% 5. Para proyectos con varios capítulos, crea un
%    anexos.tex por capítulo y agrégalo arriba con
%    \input{}.


% ==========================================
% REFERENCIA: Cómo incluir código fuente
% ==========================================
%
% Para proyectos pequeños, puedes poner los anexos
% directamente aquí. Para proyectos grandes, crea un
% archivo por capítulo (ver sección "Documentos grandes"
% en el README.md).
%
% A continuación se muestran las formas más comunes
% de incluir código con listings.
%
% ==========================================
%
% --- FORMA 1: Desde archivo externo (recomendada) ---
% Crea una carpeta codigos/ o programa/ en la raíz.
%
% \chapter{Código fuente}
% \label{anexo:codigo}
%
% \lstinputlisting[language=C, style=mystyle,
%   caption={Descripción del código. \propia{}},
%   label={lst:etiqueta}]{codigos/mi_archivo.c}
%
%
% --- FORMA 2: Con verificación de existencia ---
% Si el archivo puede no estar presente, usa \IfFileExists
% para evitar errores de compilación.
%
% \IfFileExists{programa/mi_programa.c}{
%   \lstinputlisting[language=C, style=mystyle,
%     caption={Código completo. \propia{}},
%     label={lst:codigo}]{programa/mi_programa.c}
% }{
%   \noindent\textbf{Nota:} El archivo aún no está disponible.
% }
%
%
% --- FORMA 3: Código inline (fragmentos pequeños) ---
% \begin{lstlisting}[language=Python, style=mystyle,
%   caption={Fragmento de ejemplo},
%   label={lst:ejemplo}]
% def fibonacci(n):
%     if n <= 1:
%         return n
%     return fibonacci(n - 1) + fibonacci(n - 2)
% \end{lstlisting}
%
%
% --- FORMA 4: Makefile ---
% \lstinputlisting[language=make, style=mystyle,
%   caption={Makefile del proyecto. \propia{}},
%   label={lst:makefile}]{codigos/Makefile}
%
%
% --- FORMA 5: Archivos de cabecera (.h) ---
% \lstinputlisting[language=C, style=mystyle,
%   caption={Archivo de cabecera. \propia{}},
%   label={lst:header}]{programa/common.h}
%
%
% ==========================================
% BUENAS PRÁCTICAS
% ==========================================
%
% 1. Usa \lstinputlisting para código que compila.
%    Así el PDF siempre refleja el código real.
%
% 2. Usa \IfFileExists para que la compilación no
%    falle si falta un archivo.
%
% 3. Pon siempre caption y label para referenciar
%    los listados desde el texto:
%    "En el Listado~\ref{lst:etiqueta} se muestra..."
%
% 4. Usa \propia{} al final del caption para agregar
%    "Fuente: Elaboración propia."
%
% 5. Para proyectos con varios capítulos, crea un
%    anexos.tex por capítulo y agrégalo arriba con
%    \input{}.

% % ==========================================
% ANEXOS — ARCHIVO AGREGADOR
% ==========================================
%
% Este archivo es el punto de entrada de todos los anexos.
% No pongas \chapter{} aquí directamente (a menos que sea
% un anexo general). En su lugar, crea un anexo por capítulo
% y agrégalo con \input{}.
%
% TIP: Cada capítulo puede tener su propio anexo en:
%      capitulos/capitulo_X/anexos.tex
%
% TIP: Los anexos se numeran con letras (A, B, C...)
%      automáticamente según el orden de los \input{}.
% ==========================================

\appendix
\renewcommand{\appendixname}{Anexo}
\renewcommand{\appendixpagename}{Anexos}
\renewcommand{\appendixtocname}{Anexos}
\addappheadtotoc
\appendixpage

% --- Agrega aquí los anexos de cada capítulo ---
% Ejemplo:
% % ==========================================
% ANEXOS — ARCHIVO AGREGADOR
% ==========================================
%
% Este archivo es el punto de entrada de todos los anexos.
% No pongas \chapter{} aquí directamente (a menos que sea
% un anexo general). En su lugar, crea un anexo por capítulo
% y agrégalo con \input{}.
%
% TIP: Cada capítulo puede tener su propio anexo en:
%      capitulos/capitulo_X/anexos.tex
%
% TIP: Los anexos se numeran con letras (A, B, C...)
%      automáticamente según el orden de los \input{}.
% ==========================================

\appendix
\renewcommand{\appendixname}{Anexo}
\renewcommand{\appendixpagename}{Anexos}
\renewcommand{\appendixtocname}{Anexos}
\addappheadtotoc
\appendixpage

% --- Agrega aquí los anexos de cada capítulo ---
% Ejemplo:
% \input{capitulos/capitulo_1/anexos}
% \input{capitulos/capitulo_2/anexos}

% ==========================================
% REFERENCIA: Cómo incluir código fuente
% ==========================================
%
% Para proyectos pequeños, puedes poner los anexos
% directamente aquí. Para proyectos grandes, crea un
% archivo por capítulo (ver sección "Documentos grandes"
% en el README.md).
%
% A continuación se muestran las formas más comunes
% de incluir código con listings.
%
% ==========================================
%
% --- FORMA 1: Desde archivo externo (recomendada) ---
% Crea una carpeta codigos/ o programa/ en la raíz.
%
% \chapter{Código fuente}
% \label{anexo:codigo}
%
% \lstinputlisting[language=C, style=mystyle,
%   caption={Descripción del código. \propia{}},
%   label={lst:etiqueta}]{codigos/mi_archivo.c}
%
%
% --- FORMA 2: Con verificación de existencia ---
% Si el archivo puede no estar presente, usa \IfFileExists
% para evitar errores de compilación.
%
% \IfFileExists{programa/mi_programa.c}{
%   \lstinputlisting[language=C, style=mystyle,
%     caption={Código completo. \propia{}},
%     label={lst:codigo}]{programa/mi_programa.c}
% }{
%   \noindent\textbf{Nota:} El archivo aún no está disponible.
% }
%
%
% --- FORMA 3: Código inline (fragmentos pequeños) ---
% \begin{lstlisting}[language=Python, style=mystyle,
%   caption={Fragmento de ejemplo},
%   label={lst:ejemplo}]
% def fibonacci(n):
%     if n <= 1:
%         return n
%     return fibonacci(n - 1) + fibonacci(n - 2)
% \end{lstlisting}
%
%
% --- FORMA 4: Makefile ---
% \lstinputlisting[language=make, style=mystyle,
%   caption={Makefile del proyecto. \propia{}},
%   label={lst:makefile}]{codigos/Makefile}
%
%
% --- FORMA 5: Archivos de cabecera (.h) ---
% \lstinputlisting[language=C, style=mystyle,
%   caption={Archivo de cabecera. \propia{}},
%   label={lst:header}]{programa/common.h}
%
%
% ==========================================
% BUENAS PRÁCTICAS
% ==========================================
%
% 1. Usa \lstinputlisting para código que compila.
%    Así el PDF siempre refleja el código real.
%
% 2. Usa \IfFileExists para que la compilación no
%    falle si falta un archivo.
%
% 3. Pon siempre caption y label para referenciar
%    los listados desde el texto:
%    "En el Listado~\ref{lst:etiqueta} se muestra..."
%
% 4. Usa \propia{} al final del caption para agregar
%    "Fuente: Elaboración propia."
%
% 5. Para proyectos con varios capítulos, crea un
%    anexos.tex por capítulo y agrégalo arriba con
%    \input{}.

% % ==========================================
% ANEXOS — ARCHIVO AGREGADOR
% ==========================================
%
% Este archivo es el punto de entrada de todos los anexos.
% No pongas \chapter{} aquí directamente (a menos que sea
% un anexo general). En su lugar, crea un anexo por capítulo
% y agrégalo con \input{}.
%
% TIP: Cada capítulo puede tener su propio anexo en:
%      capitulos/capitulo_X/anexos.tex
%
% TIP: Los anexos se numeran con letras (A, B, C...)
%      automáticamente según el orden de los \input{}.
% ==========================================

\appendix
\renewcommand{\appendixname}{Anexo}
\renewcommand{\appendixpagename}{Anexos}
\renewcommand{\appendixtocname}{Anexos}
\addappheadtotoc
\appendixpage

% --- Agrega aquí los anexos de cada capítulo ---
% Ejemplo:
% \input{capitulos/capitulo_1/anexos}
% \input{capitulos/capitulo_2/anexos}

% ==========================================
% REFERENCIA: Cómo incluir código fuente
% ==========================================
%
% Para proyectos pequeños, puedes poner los anexos
% directamente aquí. Para proyectos grandes, crea un
% archivo por capítulo (ver sección "Documentos grandes"
% en el README.md).
%
% A continuación se muestran las formas más comunes
% de incluir código con listings.
%
% ==========================================
%
% --- FORMA 1: Desde archivo externo (recomendada) ---
% Crea una carpeta codigos/ o programa/ en la raíz.
%
% \chapter{Código fuente}
% \label{anexo:codigo}
%
% \lstinputlisting[language=C, style=mystyle,
%   caption={Descripción del código. \propia{}},
%   label={lst:etiqueta}]{codigos/mi_archivo.c}
%
%
% --- FORMA 2: Con verificación de existencia ---
% Si el archivo puede no estar presente, usa \IfFileExists
% para evitar errores de compilación.
%
% \IfFileExists{programa/mi_programa.c}{
%   \lstinputlisting[language=C, style=mystyle,
%     caption={Código completo. \propia{}},
%     label={lst:codigo}]{programa/mi_programa.c}
% }{
%   \noindent\textbf{Nota:} El archivo aún no está disponible.
% }
%
%
% --- FORMA 3: Código inline (fragmentos pequeños) ---
% \begin{lstlisting}[language=Python, style=mystyle,
%   caption={Fragmento de ejemplo},
%   label={lst:ejemplo}]
% def fibonacci(n):
%     if n <= 1:
%         return n
%     return fibonacci(n - 1) + fibonacci(n - 2)
% \end{lstlisting}
%
%
% --- FORMA 4: Makefile ---
% \lstinputlisting[language=make, style=mystyle,
%   caption={Makefile del proyecto. \propia{}},
%   label={lst:makefile}]{codigos/Makefile}
%
%
% --- FORMA 5: Archivos de cabecera (.h) ---
% \lstinputlisting[language=C, style=mystyle,
%   caption={Archivo de cabecera. \propia{}},
%   label={lst:header}]{programa/common.h}
%
%
% ==========================================
% BUENAS PRÁCTICAS
% ==========================================
%
% 1. Usa \lstinputlisting para código que compila.
%    Así el PDF siempre refleja el código real.
%
% 2. Usa \IfFileExists para que la compilación no
%    falle si falta un archivo.
%
% 3. Pon siempre caption y label para referenciar
%    los listados desde el texto:
%    "En el Listado~\ref{lst:etiqueta} se muestra..."
%
% 4. Usa \propia{} al final del caption para agregar
%    "Fuente: Elaboración propia."
%
% 5. Para proyectos con varios capítulos, crea un
%    anexos.tex por capítulo y agrégalo arriba con
%    \input{}.


% ==========================================
% REFERENCIA: Cómo incluir código fuente
% ==========================================
%
% Para proyectos pequeños, puedes poner los anexos
% directamente aquí. Para proyectos grandes, crea un
% archivo por capítulo (ver sección "Documentos grandes"
% en el README.md).
%
% A continuación se muestran las formas más comunes
% de incluir código con listings.
%
% ==========================================
%
% --- FORMA 1: Desde archivo externo (recomendada) ---
% Crea una carpeta codigos/ o programa/ en la raíz.
%
% \chapter{Código fuente}
% \label{anexo:codigo}
%
% \lstinputlisting[language=C, style=mystyle,
%   caption={Descripción del código. \propia{}},
%   label={lst:etiqueta}]{codigos/mi_archivo.c}
%
%
% --- FORMA 2: Con verificación de existencia ---
% Si el archivo puede no estar presente, usa \IfFileExists
% para evitar errores de compilación.
%
% \IfFileExists{programa/mi_programa.c}{
%   \lstinputlisting[language=C, style=mystyle,
%     caption={Código completo. \propia{}},
%     label={lst:codigo}]{programa/mi_programa.c}
% }{
%   \noindent\textbf{Nota:} El archivo aún no está disponible.
% }
%
%
% --- FORMA 3: Código inline (fragmentos pequeños) ---
% \begin{lstlisting}[language=Python, style=mystyle,
%   caption={Fragmento de ejemplo},
%   label={lst:ejemplo}]
% def fibonacci(n):
%     if n <= 1:
%         return n
%     return fibonacci(n - 1) + fibonacci(n - 2)
% \end{lstlisting}
%
%
% --- FORMA 4: Makefile ---
% \lstinputlisting[language=make, style=mystyle,
%   caption={Makefile del proyecto. \propia{}},
%   label={lst:makefile}]{codigos/Makefile}
%
%
% --- FORMA 5: Archivos de cabecera (.h) ---
% \lstinputlisting[language=C, style=mystyle,
%   caption={Archivo de cabecera. \propia{}},
%   label={lst:header}]{programa/common.h}
%
%
% ==========================================
% BUENAS PRÁCTICAS
% ==========================================
%
% 1. Usa \lstinputlisting para código que compila.
%    Así el PDF siempre refleja el código real.
%
% 2. Usa \IfFileExists para que la compilación no
%    falle si falta un archivo.
%
% 3. Pon siempre caption y label para referenciar
%    los listados desde el texto:
%    "En el Listado~\ref{lst:etiqueta} se muestra..."
%
% 4. Usa \propia{} al final del caption para agregar
%    "Fuente: Elaboración propia."
%
% 5. Para proyectos con varios capítulos, crea un
%    anexos.tex por capítulo y agrégalo arriba con
%    \input{}.


% ==========================================
% REFERENCIA: Cómo incluir código fuente
% ==========================================
%
% Para proyectos pequeños, puedes poner los anexos
% directamente aquí. Para proyectos grandes, crea un
% archivo por capítulo (ver sección "Documentos grandes"
% en el README.md).
%
% A continuación se muestran las formas más comunes
% de incluir código con listings.
%
% ==========================================
%
% --- FORMA 1: Desde archivo externo (recomendada) ---
% Crea una carpeta codigos/ o programa/ en la raíz.
%
% \chapter{Código fuente}
% \label{anexo:codigo}
%
% \lstinputlisting[language=C, style=mystyle,
%   caption={Descripción del código. \propia{}},
%   label={lst:etiqueta}]{codigos/mi_archivo.c}
%
%
% --- FORMA 2: Con verificación de existencia ---
% Si el archivo puede no estar presente, usa \IfFileExists
% para evitar errores de compilación.
%
% \IfFileExists{programa/mi_programa.c}{
%   \lstinputlisting[language=C, style=mystyle,
%     caption={Código completo. \propia{}},
%     label={lst:codigo}]{programa/mi_programa.c}
% }{
%   \noindent\textbf{Nota:} El archivo aún no está disponible.
% }
%
%
% --- FORMA 3: Código inline (fragmentos pequeños) ---
% \begin{lstlisting}[language=Python, style=mystyle,
%   caption={Fragmento de ejemplo},
%   label={lst:ejemplo}]
% def fibonacci(n):
%     if n <= 1:
%         return n
%     return fibonacci(n - 1) + fibonacci(n - 2)
% \end{lstlisting}
%
%
% --- FORMA 4: Makefile ---
% \lstinputlisting[language=make, style=mystyle,
%   caption={Makefile del proyecto. \propia{}},
%   label={lst:makefile}]{codigos/Makefile}
%
%
% --- FORMA 5: Archivos de cabecera (.h) ---
% \lstinputlisting[language=C, style=mystyle,
%   caption={Archivo de cabecera. \propia{}},
%   label={lst:header}]{programa/common.h}
%
%
% ==========================================
% BUENAS PRÁCTICAS
% ==========================================
%
% 1. Usa \lstinputlisting para código que compila.
%    Así el PDF siempre refleja el código real.
%
% 2. Usa \IfFileExists para que la compilación no
%    falle si falta un archivo.
%
% 3. Pon siempre caption y label para referenciar
%    los listados desde el texto:
%    "En el Listado~\ref{lst:etiqueta} se muestra..."
%
% 4. Usa \propia{} al final del caption para agregar
%    "Fuente: Elaboración propia."
%
% 5. Para proyectos con varios capítulos, crea un
%    anexos.tex por capítulo y agrégalo arriba con
%    \input{}.

% % ==========================================
% ANEXOS — ARCHIVO AGREGADOR
% ==========================================
%
% Este archivo es el punto de entrada de todos los anexos.
% No pongas \chapter{} aquí directamente (a menos que sea
% un anexo general). En su lugar, crea un anexo por capítulo
% y agrégalo con \input{}.
%
% TIP: Cada capítulo puede tener su propio anexo en:
%      capitulos/capitulo_X/anexos.tex
%
% TIP: Los anexos se numeran con letras (A, B, C...)
%      automáticamente según el orden de los \input{}.
% ==========================================

\appendix
\renewcommand{\appendixname}{Anexo}
\renewcommand{\appendixpagename}{Anexos}
\renewcommand{\appendixtocname}{Anexos}
\addappheadtotoc
\appendixpage

% --- Agrega aquí los anexos de cada capítulo ---
% Ejemplo:
% % ==========================================
% ANEXOS — ARCHIVO AGREGADOR
% ==========================================
%
% Este archivo es el punto de entrada de todos los anexos.
% No pongas \chapter{} aquí directamente (a menos que sea
% un anexo general). En su lugar, crea un anexo por capítulo
% y agrégalo con \input{}.
%
% TIP: Cada capítulo puede tener su propio anexo en:
%      capitulos/capitulo_X/anexos.tex
%
% TIP: Los anexos se numeran con letras (A, B, C...)
%      automáticamente según el orden de los \input{}.
% ==========================================

\appendix
\renewcommand{\appendixname}{Anexo}
\renewcommand{\appendixpagename}{Anexos}
\renewcommand{\appendixtocname}{Anexos}
\addappheadtotoc
\appendixpage

% --- Agrega aquí los anexos de cada capítulo ---
% Ejemplo:
% % ==========================================
% ANEXOS — ARCHIVO AGREGADOR
% ==========================================
%
% Este archivo es el punto de entrada de todos los anexos.
% No pongas \chapter{} aquí directamente (a menos que sea
% un anexo general). En su lugar, crea un anexo por capítulo
% y agrégalo con \input{}.
%
% TIP: Cada capítulo puede tener su propio anexo en:
%      capitulos/capitulo_X/anexos.tex
%
% TIP: Los anexos se numeran con letras (A, B, C...)
%      automáticamente según el orden de los \input{}.
% ==========================================

\appendix
\renewcommand{\appendixname}{Anexo}
\renewcommand{\appendixpagename}{Anexos}
\renewcommand{\appendixtocname}{Anexos}
\addappheadtotoc
\appendixpage

% --- Agrega aquí los anexos de cada capítulo ---
% Ejemplo:
% \input{capitulos/capitulo_1/anexos}
% \input{capitulos/capitulo_2/anexos}

% ==========================================
% REFERENCIA: Cómo incluir código fuente
% ==========================================
%
% Para proyectos pequeños, puedes poner los anexos
% directamente aquí. Para proyectos grandes, crea un
% archivo por capítulo (ver sección "Documentos grandes"
% en el README.md).
%
% A continuación se muestran las formas más comunes
% de incluir código con listings.
%
% ==========================================
%
% --- FORMA 1: Desde archivo externo (recomendada) ---
% Crea una carpeta codigos/ o programa/ en la raíz.
%
% \chapter{Código fuente}
% \label{anexo:codigo}
%
% \lstinputlisting[language=C, style=mystyle,
%   caption={Descripción del código. \propia{}},
%   label={lst:etiqueta}]{codigos/mi_archivo.c}
%
%
% --- FORMA 2: Con verificación de existencia ---
% Si el archivo puede no estar presente, usa \IfFileExists
% para evitar errores de compilación.
%
% \IfFileExists{programa/mi_programa.c}{
%   \lstinputlisting[language=C, style=mystyle,
%     caption={Código completo. \propia{}},
%     label={lst:codigo}]{programa/mi_programa.c}
% }{
%   \noindent\textbf{Nota:} El archivo aún no está disponible.
% }
%
%
% --- FORMA 3: Código inline (fragmentos pequeños) ---
% \begin{lstlisting}[language=Python, style=mystyle,
%   caption={Fragmento de ejemplo},
%   label={lst:ejemplo}]
% def fibonacci(n):
%     if n <= 1:
%         return n
%     return fibonacci(n - 1) + fibonacci(n - 2)
% \end{lstlisting}
%
%
% --- FORMA 4: Makefile ---
% \lstinputlisting[language=make, style=mystyle,
%   caption={Makefile del proyecto. \propia{}},
%   label={lst:makefile}]{codigos/Makefile}
%
%
% --- FORMA 5: Archivos de cabecera (.h) ---
% \lstinputlisting[language=C, style=mystyle,
%   caption={Archivo de cabecera. \propia{}},
%   label={lst:header}]{programa/common.h}
%
%
% ==========================================
% BUENAS PRÁCTICAS
% ==========================================
%
% 1. Usa \lstinputlisting para código que compila.
%    Así el PDF siempre refleja el código real.
%
% 2. Usa \IfFileExists para que la compilación no
%    falle si falta un archivo.
%
% 3. Pon siempre caption y label para referenciar
%    los listados desde el texto:
%    "En el Listado~\ref{lst:etiqueta} se muestra..."
%
% 4. Usa \propia{} al final del caption para agregar
%    "Fuente: Elaboración propia."
%
% 5. Para proyectos con varios capítulos, crea un
%    anexos.tex por capítulo y agrégalo arriba con
%    \input{}.

% % ==========================================
% ANEXOS — ARCHIVO AGREGADOR
% ==========================================
%
% Este archivo es el punto de entrada de todos los anexos.
% No pongas \chapter{} aquí directamente (a menos que sea
% un anexo general). En su lugar, crea un anexo por capítulo
% y agrégalo con \input{}.
%
% TIP: Cada capítulo puede tener su propio anexo en:
%      capitulos/capitulo_X/anexos.tex
%
% TIP: Los anexos se numeran con letras (A, B, C...)
%      automáticamente según el orden de los \input{}.
% ==========================================

\appendix
\renewcommand{\appendixname}{Anexo}
\renewcommand{\appendixpagename}{Anexos}
\renewcommand{\appendixtocname}{Anexos}
\addappheadtotoc
\appendixpage

% --- Agrega aquí los anexos de cada capítulo ---
% Ejemplo:
% \input{capitulos/capitulo_1/anexos}
% \input{capitulos/capitulo_2/anexos}

% ==========================================
% REFERENCIA: Cómo incluir código fuente
% ==========================================
%
% Para proyectos pequeños, puedes poner los anexos
% directamente aquí. Para proyectos grandes, crea un
% archivo por capítulo (ver sección "Documentos grandes"
% en el README.md).
%
% A continuación se muestran las formas más comunes
% de incluir código con listings.
%
% ==========================================
%
% --- FORMA 1: Desde archivo externo (recomendada) ---
% Crea una carpeta codigos/ o programa/ en la raíz.
%
% \chapter{Código fuente}
% \label{anexo:codigo}
%
% \lstinputlisting[language=C, style=mystyle,
%   caption={Descripción del código. \propia{}},
%   label={lst:etiqueta}]{codigos/mi_archivo.c}
%
%
% --- FORMA 2: Con verificación de existencia ---
% Si el archivo puede no estar presente, usa \IfFileExists
% para evitar errores de compilación.
%
% \IfFileExists{programa/mi_programa.c}{
%   \lstinputlisting[language=C, style=mystyle,
%     caption={Código completo. \propia{}},
%     label={lst:codigo}]{programa/mi_programa.c}
% }{
%   \noindent\textbf{Nota:} El archivo aún no está disponible.
% }
%
%
% --- FORMA 3: Código inline (fragmentos pequeños) ---
% \begin{lstlisting}[language=Python, style=mystyle,
%   caption={Fragmento de ejemplo},
%   label={lst:ejemplo}]
% def fibonacci(n):
%     if n <= 1:
%         return n
%     return fibonacci(n - 1) + fibonacci(n - 2)
% \end{lstlisting}
%
%
% --- FORMA 4: Makefile ---
% \lstinputlisting[language=make, style=mystyle,
%   caption={Makefile del proyecto. \propia{}},
%   label={lst:makefile}]{codigos/Makefile}
%
%
% --- FORMA 5: Archivos de cabecera (.h) ---
% \lstinputlisting[language=C, style=mystyle,
%   caption={Archivo de cabecera. \propia{}},
%   label={lst:header}]{programa/common.h}
%
%
% ==========================================
% BUENAS PRÁCTICAS
% ==========================================
%
% 1. Usa \lstinputlisting para código que compila.
%    Así el PDF siempre refleja el código real.
%
% 2. Usa \IfFileExists para que la compilación no
%    falle si falta un archivo.
%
% 3. Pon siempre caption y label para referenciar
%    los listados desde el texto:
%    "En el Listado~\ref{lst:etiqueta} se muestra..."
%
% 4. Usa \propia{} al final del caption para agregar
%    "Fuente: Elaboración propia."
%
% 5. Para proyectos con varios capítulos, crea un
%    anexos.tex por capítulo y agrégalo arriba con
%    \input{}.


% ==========================================
% REFERENCIA: Cómo incluir código fuente
% ==========================================
%
% Para proyectos pequeños, puedes poner los anexos
% directamente aquí. Para proyectos grandes, crea un
% archivo por capítulo (ver sección "Documentos grandes"
% en el README.md).
%
% A continuación se muestran las formas más comunes
% de incluir código con listings.
%
% ==========================================
%
% --- FORMA 1: Desde archivo externo (recomendada) ---
% Crea una carpeta codigos/ o programa/ en la raíz.
%
% \chapter{Código fuente}
% \label{anexo:codigo}
%
% \lstinputlisting[language=C, style=mystyle,
%   caption={Descripción del código. \propia{}},
%   label={lst:etiqueta}]{codigos/mi_archivo.c}
%
%
% --- FORMA 2: Con verificación de existencia ---
% Si el archivo puede no estar presente, usa \IfFileExists
% para evitar errores de compilación.
%
% \IfFileExists{programa/mi_programa.c}{
%   \lstinputlisting[language=C, style=mystyle,
%     caption={Código completo. \propia{}},
%     label={lst:codigo}]{programa/mi_programa.c}
% }{
%   \noindent\textbf{Nota:} El archivo aún no está disponible.
% }
%
%
% --- FORMA 3: Código inline (fragmentos pequeños) ---
% \begin{lstlisting}[language=Python, style=mystyle,
%   caption={Fragmento de ejemplo},
%   label={lst:ejemplo}]
% def fibonacci(n):
%     if n <= 1:
%         return n
%     return fibonacci(n - 1) + fibonacci(n - 2)
% \end{lstlisting}
%
%
% --- FORMA 4: Makefile ---
% \lstinputlisting[language=make, style=mystyle,
%   caption={Makefile del proyecto. \propia{}},
%   label={lst:makefile}]{codigos/Makefile}
%
%
% --- FORMA 5: Archivos de cabecera (.h) ---
% \lstinputlisting[language=C, style=mystyle,
%   caption={Archivo de cabecera. \propia{}},
%   label={lst:header}]{programa/common.h}
%
%
% ==========================================
% BUENAS PRÁCTICAS
% ==========================================
%
% 1. Usa \lstinputlisting para código que compila.
%    Así el PDF siempre refleja el código real.
%
% 2. Usa \IfFileExists para que la compilación no
%    falle si falta un archivo.
%
% 3. Pon siempre caption y label para referenciar
%    los listados desde el texto:
%    "En el Listado~\ref{lst:etiqueta} se muestra..."
%
% 4. Usa \propia{} al final del caption para agregar
%    "Fuente: Elaboración propia."
%
% 5. Para proyectos con varios capítulos, crea un
%    anexos.tex por capítulo y agrégalo arriba con
%    \input{}.

% % ==========================================
% ANEXOS — ARCHIVO AGREGADOR
% ==========================================
%
% Este archivo es el punto de entrada de todos los anexos.
% No pongas \chapter{} aquí directamente (a menos que sea
% un anexo general). En su lugar, crea un anexo por capítulo
% y agrégalo con \input{}.
%
% TIP: Cada capítulo puede tener su propio anexo en:
%      capitulos/capitulo_X/anexos.tex
%
% TIP: Los anexos se numeran con letras (A, B, C...)
%      automáticamente según el orden de los \input{}.
% ==========================================

\appendix
\renewcommand{\appendixname}{Anexo}
\renewcommand{\appendixpagename}{Anexos}
\renewcommand{\appendixtocname}{Anexos}
\addappheadtotoc
\appendixpage

% --- Agrega aquí los anexos de cada capítulo ---
% Ejemplo:
% % ==========================================
% ANEXOS — ARCHIVO AGREGADOR
% ==========================================
%
% Este archivo es el punto de entrada de todos los anexos.
% No pongas \chapter{} aquí directamente (a menos que sea
% un anexo general). En su lugar, crea un anexo por capítulo
% y agrégalo con \input{}.
%
% TIP: Cada capítulo puede tener su propio anexo en:
%      capitulos/capitulo_X/anexos.tex
%
% TIP: Los anexos se numeran con letras (A, B, C...)
%      automáticamente según el orden de los \input{}.
% ==========================================

\appendix
\renewcommand{\appendixname}{Anexo}
\renewcommand{\appendixpagename}{Anexos}
\renewcommand{\appendixtocname}{Anexos}
\addappheadtotoc
\appendixpage

% --- Agrega aquí los anexos de cada capítulo ---
% Ejemplo:
% \input{capitulos/capitulo_1/anexos}
% \input{capitulos/capitulo_2/anexos}

% ==========================================
% REFERENCIA: Cómo incluir código fuente
% ==========================================
%
% Para proyectos pequeños, puedes poner los anexos
% directamente aquí. Para proyectos grandes, crea un
% archivo por capítulo (ver sección "Documentos grandes"
% en el README.md).
%
% A continuación se muestran las formas más comunes
% de incluir código con listings.
%
% ==========================================
%
% --- FORMA 1: Desde archivo externo (recomendada) ---
% Crea una carpeta codigos/ o programa/ en la raíz.
%
% \chapter{Código fuente}
% \label{anexo:codigo}
%
% \lstinputlisting[language=C, style=mystyle,
%   caption={Descripción del código. \propia{}},
%   label={lst:etiqueta}]{codigos/mi_archivo.c}
%
%
% --- FORMA 2: Con verificación de existencia ---
% Si el archivo puede no estar presente, usa \IfFileExists
% para evitar errores de compilación.
%
% \IfFileExists{programa/mi_programa.c}{
%   \lstinputlisting[language=C, style=mystyle,
%     caption={Código completo. \propia{}},
%     label={lst:codigo}]{programa/mi_programa.c}
% }{
%   \noindent\textbf{Nota:} El archivo aún no está disponible.
% }
%
%
% --- FORMA 3: Código inline (fragmentos pequeños) ---
% \begin{lstlisting}[language=Python, style=mystyle,
%   caption={Fragmento de ejemplo},
%   label={lst:ejemplo}]
% def fibonacci(n):
%     if n <= 1:
%         return n
%     return fibonacci(n - 1) + fibonacci(n - 2)
% \end{lstlisting}
%
%
% --- FORMA 4: Makefile ---
% \lstinputlisting[language=make, style=mystyle,
%   caption={Makefile del proyecto. \propia{}},
%   label={lst:makefile}]{codigos/Makefile}
%
%
% --- FORMA 5: Archivos de cabecera (.h) ---
% \lstinputlisting[language=C, style=mystyle,
%   caption={Archivo de cabecera. \propia{}},
%   label={lst:header}]{programa/common.h}
%
%
% ==========================================
% BUENAS PRÁCTICAS
% ==========================================
%
% 1. Usa \lstinputlisting para código que compila.
%    Así el PDF siempre refleja el código real.
%
% 2. Usa \IfFileExists para que la compilación no
%    falle si falta un archivo.
%
% 3. Pon siempre caption y label para referenciar
%    los listados desde el texto:
%    "En el Listado~\ref{lst:etiqueta} se muestra..."
%
% 4. Usa \propia{} al final del caption para agregar
%    "Fuente: Elaboración propia."
%
% 5. Para proyectos con varios capítulos, crea un
%    anexos.tex por capítulo y agrégalo arriba con
%    \input{}.

% % ==========================================
% ANEXOS — ARCHIVO AGREGADOR
% ==========================================
%
% Este archivo es el punto de entrada de todos los anexos.
% No pongas \chapter{} aquí directamente (a menos que sea
% un anexo general). En su lugar, crea un anexo por capítulo
% y agrégalo con \input{}.
%
% TIP: Cada capítulo puede tener su propio anexo en:
%      capitulos/capitulo_X/anexos.tex
%
% TIP: Los anexos se numeran con letras (A, B, C...)
%      automáticamente según el orden de los \input{}.
% ==========================================

\appendix
\renewcommand{\appendixname}{Anexo}
\renewcommand{\appendixpagename}{Anexos}
\renewcommand{\appendixtocname}{Anexos}
\addappheadtotoc
\appendixpage

% --- Agrega aquí los anexos de cada capítulo ---
% Ejemplo:
% \input{capitulos/capitulo_1/anexos}
% \input{capitulos/capitulo_2/anexos}

% ==========================================
% REFERENCIA: Cómo incluir código fuente
% ==========================================
%
% Para proyectos pequeños, puedes poner los anexos
% directamente aquí. Para proyectos grandes, crea un
% archivo por capítulo (ver sección "Documentos grandes"
% en el README.md).
%
% A continuación se muestran las formas más comunes
% de incluir código con listings.
%
% ==========================================
%
% --- FORMA 1: Desde archivo externo (recomendada) ---
% Crea una carpeta codigos/ o programa/ en la raíz.
%
% \chapter{Código fuente}
% \label{anexo:codigo}
%
% \lstinputlisting[language=C, style=mystyle,
%   caption={Descripción del código. \propia{}},
%   label={lst:etiqueta}]{codigos/mi_archivo.c}
%
%
% --- FORMA 2: Con verificación de existencia ---
% Si el archivo puede no estar presente, usa \IfFileExists
% para evitar errores de compilación.
%
% \IfFileExists{programa/mi_programa.c}{
%   \lstinputlisting[language=C, style=mystyle,
%     caption={Código completo. \propia{}},
%     label={lst:codigo}]{programa/mi_programa.c}
% }{
%   \noindent\textbf{Nota:} El archivo aún no está disponible.
% }
%
%
% --- FORMA 3: Código inline (fragmentos pequeños) ---
% \begin{lstlisting}[language=Python, style=mystyle,
%   caption={Fragmento de ejemplo},
%   label={lst:ejemplo}]
% def fibonacci(n):
%     if n <= 1:
%         return n
%     return fibonacci(n - 1) + fibonacci(n - 2)
% \end{lstlisting}
%
%
% --- FORMA 4: Makefile ---
% \lstinputlisting[language=make, style=mystyle,
%   caption={Makefile del proyecto. \propia{}},
%   label={lst:makefile}]{codigos/Makefile}
%
%
% --- FORMA 5: Archivos de cabecera (.h) ---
% \lstinputlisting[language=C, style=mystyle,
%   caption={Archivo de cabecera. \propia{}},
%   label={lst:header}]{programa/common.h}
%
%
% ==========================================
% BUENAS PRÁCTICAS
% ==========================================
%
% 1. Usa \lstinputlisting para código que compila.
%    Así el PDF siempre refleja el código real.
%
% 2. Usa \IfFileExists para que la compilación no
%    falle si falta un archivo.
%
% 3. Pon siempre caption y label para referenciar
%    los listados desde el texto:
%    "En el Listado~\ref{lst:etiqueta} se muestra..."
%
% 4. Usa \propia{} al final del caption para agregar
%    "Fuente: Elaboración propia."
%
% 5. Para proyectos con varios capítulos, crea un
%    anexos.tex por capítulo y agrégalo arriba con
%    \input{}.


% ==========================================
% REFERENCIA: Cómo incluir código fuente
% ==========================================
%
% Para proyectos pequeños, puedes poner los anexos
% directamente aquí. Para proyectos grandes, crea un
% archivo por capítulo (ver sección "Documentos grandes"
% en el README.md).
%
% A continuación se muestran las formas más comunes
% de incluir código con listings.
%
% ==========================================
%
% --- FORMA 1: Desde archivo externo (recomendada) ---
% Crea una carpeta codigos/ o programa/ en la raíz.
%
% \chapter{Código fuente}
% \label{anexo:codigo}
%
% \lstinputlisting[language=C, style=mystyle,
%   caption={Descripción del código. \propia{}},
%   label={lst:etiqueta}]{codigos/mi_archivo.c}
%
%
% --- FORMA 2: Con verificación de existencia ---
% Si el archivo puede no estar presente, usa \IfFileExists
% para evitar errores de compilación.
%
% \IfFileExists{programa/mi_programa.c}{
%   \lstinputlisting[language=C, style=mystyle,
%     caption={Código completo. \propia{}},
%     label={lst:codigo}]{programa/mi_programa.c}
% }{
%   \noindent\textbf{Nota:} El archivo aún no está disponible.
% }
%
%
% --- FORMA 3: Código inline (fragmentos pequeños) ---
% \begin{lstlisting}[language=Python, style=mystyle,
%   caption={Fragmento de ejemplo},
%   label={lst:ejemplo}]
% def fibonacci(n):
%     if n <= 1:
%         return n
%     return fibonacci(n - 1) + fibonacci(n - 2)
% \end{lstlisting}
%
%
% --- FORMA 4: Makefile ---
% \lstinputlisting[language=make, style=mystyle,
%   caption={Makefile del proyecto. \propia{}},
%   label={lst:makefile}]{codigos/Makefile}
%
%
% --- FORMA 5: Archivos de cabecera (.h) ---
% \lstinputlisting[language=C, style=mystyle,
%   caption={Archivo de cabecera. \propia{}},
%   label={lst:header}]{programa/common.h}
%
%
% ==========================================
% BUENAS PRÁCTICAS
% ==========================================
%
% 1. Usa \lstinputlisting para código que compila.
%    Así el PDF siempre refleja el código real.
%
% 2. Usa \IfFileExists para que la compilación no
%    falle si falta un archivo.
%
% 3. Pon siempre caption y label para referenciar
%    los listados desde el texto:
%    "En el Listado~\ref{lst:etiqueta} se muestra..."
%
% 4. Usa \propia{} al final del caption para agregar
%    "Fuente: Elaboración propia."
%
% 5. Para proyectos con varios capítulos, crea un
%    anexos.tex por capítulo y agrégalo arriba con
%    \input{}.


% ==========================================
% REFERENCIA: Cómo incluir código fuente
% ==========================================
%
% Para proyectos pequeños, puedes poner los anexos
% directamente aquí. Para proyectos grandes, crea un
% archivo por capítulo (ver sección "Documentos grandes"
% en el README.md).
%
% A continuación se muestran las formas más comunes
% de incluir código con listings.
%
% ==========================================
%
% --- FORMA 1: Desde archivo externo (recomendada) ---
% Crea una carpeta codigos/ o programa/ en la raíz.
%
% \chapter{Código fuente}
% \label{anexo:codigo}
%
% \lstinputlisting[language=C, style=mystyle,
%   caption={Descripción del código. \propia{}},
%   label={lst:etiqueta}]{codigos/mi_archivo.c}
%
%
% --- FORMA 2: Con verificación de existencia ---
% Si el archivo puede no estar presente, usa \IfFileExists
% para evitar errores de compilación.
%
% \IfFileExists{programa/mi_programa.c}{
%   \lstinputlisting[language=C, style=mystyle,
%     caption={Código completo. \propia{}},
%     label={lst:codigo}]{programa/mi_programa.c}
% }{
%   \noindent\textbf{Nota:} El archivo aún no está disponible.
% }
%
%
% --- FORMA 3: Código inline (fragmentos pequeños) ---
% \begin{lstlisting}[language=Python, style=mystyle,
%   caption={Fragmento de ejemplo},
%   label={lst:ejemplo}]
% def fibonacci(n):
%     if n <= 1:
%         return n
%     return fibonacci(n - 1) + fibonacci(n - 2)
% \end{lstlisting}
%
%
% --- FORMA 4: Makefile ---
% \lstinputlisting[language=make, style=mystyle,
%   caption={Makefile del proyecto. \propia{}},
%   label={lst:makefile}]{codigos/Makefile}
%
%
% --- FORMA 5: Archivos de cabecera (.h) ---
% \lstinputlisting[language=C, style=mystyle,
%   caption={Archivo de cabecera. \propia{}},
%   label={lst:header}]{programa/common.h}
%
%
% ==========================================
% BUENAS PRÁCTICAS
% ==========================================
%
% 1. Usa \lstinputlisting para código que compila.
%    Así el PDF siempre refleja el código real.
%
% 2. Usa \IfFileExists para que la compilación no
%    falle si falta un archivo.
%
% 3. Pon siempre caption y label para referenciar
%    los listados desde el texto:
%    "En el Listado~\ref{lst:etiqueta} se muestra..."
%
% 4. Usa \propia{} al final del caption para agregar
%    "Fuente: Elaboración propia."
%
% 5. Para proyectos con varios capítulos, crea un
%    anexos.tex por capítulo y agrégalo arriba con
%    \input{}.


% ==========================================
% REFERENCIA: Cómo incluir código fuente
% ==========================================
%
% Para proyectos pequeños, puedes poner los anexos
% directamente aquí. Para proyectos grandes, crea un
% archivo por capítulo (ver sección "Documentos grandes"
% en el README.md).
%
% A continuación se muestran las formas más comunes
% de incluir código con listings.
%
% ==========================================
%
% --- FORMA 1: Desde archivo externo (recomendada) ---
% Crea una carpeta codigos/ o programa/ en la raíz.
%
% \chapter{Código fuente}
% \label{anexo:codigo}
%
% \lstinputlisting[language=C, style=mystyle,
%   caption={Descripción del código. \propia{}},
%   label={lst:etiqueta}]{codigos/mi_archivo.c}
%
%
% --- FORMA 2: Con verificación de existencia ---
% Si el archivo puede no estar presente, usa \IfFileExists
% para evitar errores de compilación.
%
% \IfFileExists{programa/mi_programa.c}{
%   \lstinputlisting[language=C, style=mystyle,
%     caption={Código completo. \propia{}},
%     label={lst:codigo}]{programa/mi_programa.c}
% }{
%   \noindent\textbf{Nota:} El archivo aún no está disponible.
% }
%
%
% --- FORMA 3: Código inline (fragmentos pequeños) ---
% \begin{lstlisting}[language=Python, style=mystyle,
%   caption={Fragmento de ejemplo},
%   label={lst:ejemplo}]
% def fibonacci(n):
%     if n <= 1:
%         return n
%     return fibonacci(n - 1) + fibonacci(n - 2)
% \end{lstlisting}
%
%
% --- FORMA 4: Makefile ---
% \lstinputlisting[language=make, style=mystyle,
%   caption={Makefile del proyecto. \propia{}},
%   label={lst:makefile}]{codigos/Makefile}
%
%
% --- FORMA 5: Archivos de cabecera (.h) ---
% \lstinputlisting[language=C, style=mystyle,
%   caption={Archivo de cabecera. \propia{}},
%   label={lst:header}]{programa/common.h}
%
%
% ==========================================
% BUENAS PRÁCTICAS
% ==========================================
%
% 1. Usa \lstinputlisting para código que compila.
%    Así el PDF siempre refleja el código real.
%
% 2. Usa \IfFileExists para que la compilación no
%    falle si falta un archivo.
%
% 3. Pon siempre caption y label para referenciar
%    los listados desde el texto:
%    "En el Listado~\ref{lst:etiqueta} se muestra..."
%
% 4. Usa \propia{} al final del caption para agregar
%    "Fuente: Elaboración propia."
%
% 5. Para proyectos con varios capítulos, crea un
%    anexos.tex por capítulo y agrégalo arriba con
%    \input{}.
